\section{Constructing a nonradial probability path and its correction}
\label{sec:nonradial-construction}

The acceleration condition can be imposed before solving a nonlinear
Eulerian equation. Let $\rho$ be a reference law, $b$ a field preserving
$\rho$, and $F_t$ smooth diffeomorphisms. If
\begin{equation}
 \begin{aligned}
  \partial_{tt}F_t(z)&=D F_t(z)b(z),\\
  p_t&=(F_t)_\#\rho,\\
  u_t(F_tz)&=\partial_tF_t(z).
 \end{aligned}
 \label{eq:reference-acceleration}
\end{equation}
then $\acc_t(F_tz)=DF_t(z)b(z)$. For every compactly supported smooth test
$\phi$, the change of variables gives
$\int\nabla\phi\cdot\acc_t\,\dd p_t
=\int\nabla(\phi\circ F_t)\cdot b\,\dd\rho=0$.
Thus the linear equation for $F_t$ in
Eq.~\eqref{eq:reference-acceleration} constructs an admissible acceleration.
A general nonlinear pushforward does not have this property: the displayed
second-derivative identity is essential.

\subsection{An explicit globally regular family}

Identify $\R^2$ with $\mathbb C$, so $Jz=iz$. Take
$\rho=\mathcal N(0,I_2)$, $a,T>0$, and $\epsilon>0$. Define
\begin{align}
 f(z)&=z^2e^{-|z|^2},\qquad
 A_t=e^{(1+i)at},\qquad B_t=\epsilon e^{(1+i)a\sqrt2t},\nonumber\\
 F_t(z)&=A_tz+B_tf(z),\qquad
 q=\frac{4\epsilon}{e}e^{a(\sqrt2-1)T}<1.
 \label{eq:harmonic-map}
\end{align}
The source law is $\rho$; the initial target law is $p_0=(F_0)_\#\rho$,
which is already non-Gaussian. Writing
$k_t(F_tz)=DF_t(z)Jz$, set
\begin{equation}
 v_t(F_tz)=\partial_tF_t(z)-a k_t(F_tz),\qquad
 u_{\omega,t}=v_t+\omega k_t.
 \label{eq:harmonic-gauge}
\end{equation}
Here $v$ is a base generator, not a claimed minimum-action field.

\begin{theorem}[A constructive nonradial family]
\label{thm:harmonic-construction}
Under Eq.~\eqref{eq:harmonic-map}, each $F_t$ is a global smooth
bi-Lipschitz diffeomorphism, uniformly on $[0,T]$. The laws
$p_t=(F_t)_\#\rho$ have smooth positive densities on $\R^2$ and finite
moments of every order. For every fixed $\omega\in\R$,
$u_\omega$ generates exactly this same path, with trajectories
\[
 X_t=F_t\bigl(e^{i(\omega-a)t}Z\bigr),\qquad Z\sim\rho.
\]
The choice $\omega=a$ satisfies the acceleration condition, and all
smoothness, growth, global Lipschitz, and moment assumptions of
Theorem~\ref{thm:distribution_order}. Consequently its Euler distributions
have error $O(H^2)$ on every finite grid. For $\epsilon>0$, no $p_t$ is
an invertible affine image of a radial law or of a product of two
independent real laws.
\end{theorem}

The key identity is
\[
 \partial_{tt}F_t=2a^2\partial_\theta F_t
                 =2a^2DF_tJz.
\]
It follows because the angular harmonics $z$ and $f(z)$ have degrees one
and two, and their time exponents have squares $2ia^2$ and $4ia^2$.
Also $\|Df\|_\infty\le4/e$, so $F_t^{-1}(x)$ is the unique limit of
the contraction
$z_{m+1}=A_t^{-1}(x-B_tf(z_m))$, with contraction factor at most $q$.
Thus the density
$p_t(x)=\rho(F_t^{-1}x)/\det DF_t(F_t^{-1}x)$ and the velocity are
computable without a PDE solve. The inverse iteration adds real
computational cost; the numerical comparison includes it.

This is a family, not just one isolated deformation. Replace $B_tf$ by
$\sum_{n=2}^m\epsilon_n e^{(1+i)a\sqrt n t}f_n(z)$, where
$f_n(e^{i\theta}z)=e^{in\theta}f_n(z)$ and the $f_n$ are smooth rapidly
decaying functions. The same construction works whenever
\[
 \sum_{n=2}^m|\epsilon_n|e^{a(\sqrt n-1)T}\|Df_n\|_\infty<1.
\]
Arbitrary smooth radial profiles times $z^n$ give such harmonics.
This constructs a specified class of paths. It does not solve the
inverse problem for an arbitrary externally supplied density path.
Appendix~\ref{app:harmonic-proofs} proves all assertions, including the
nonproduct certificate.

\begin{corollary}[A Gaussian starting law]
\label{cor:gaussian-start}
Replace $B_t$ by $\epsilon\sinh((1+i)a\sqrt2t)$ and assume
$(4/e)\epsilon\sqrt2\sinh(a\sqrt2T)<1$.
Then $F_0$ is the identity, so $p_0=\mathcal N(0,I_2)$, and the
acceleration-preserving construction and its $O(H^2)$ bound still hold.
For every $t>0$ the target law is neither affine-radial nor affine-product.
\end{corollary}

This follows from $B_t''=4ia^2B_t$ and
$|\sinh((1+i)s)|^2=\sinh^2s+\sin^2s\le2\sinh^2s$.
The numerical study below uses the exponential-coefficient family of
Eq.~\eqref{eq:harmonic-map}; the Gaussian-start variant is an additional
proved construction, not a separately reported empirical result.

\subsection{Discovering the correction by a finite global search}

For the two-mode family let $F_{1,t}=A_tz$, $F_{2,t}=B_tf(z)$, and
\[
 \zeta_n=a\sqrt n+i\{a\sqrt n+n(\omega-a)\}.
\]
Then $\acc_\omega(t,F_tz)=\sum_{n=1}^2\zeta_n^2F_{n,t}(z)$.
Every $c k_t$ preserves $p_t$, so the following is an upper certificate
for the projected acceleration residual:
\begin{align}
 R_t(\omega,c)^2
 &=\|\acc_\omega-c k_t\|_{L^2(p_t)}^2\nonumber\\
 &=\sum_{n=1}^2 M_n(t)|\zeta_n^2-inc|^2,\label{eq:harmonic-residual}\\
 M_1(t)&=2e^{2at},\qquad
 M_2(t)=\frac8{125}\epsilon^2e^{2a\sqrt2t}.\nonumber
\end{align}
Angular orthogonality makes this formula exact for the continuum law.
It avoids numerical differentiation and a weighted Poisson solve.

\begin{proposition}[Global correction recovery in the stated family]
\label{prop:harmonic-search}
At any fixed $t$, minimizing Eq.~\eqref{eq:harmonic-residual} over
$(\omega,c)\in\R^2$ has the unique zero $(a,2a^2)$.
For a fixed $\omega$ the unique minimizing coefficient is
\[
 c_*(\omega)=
 \frac{\sum_{n=1}^2nM_n\operatorname{Im}(\zeta_n^2)}
      {\sum_{n=1}^2n^2M_n}.
\]
After substitution, $R_t(\omega,c_*(\omega))^2$ is a coercive quartic
polynomial. Its global minimum is found by evaluating its real cubic
stationary points; on a bounded interval include the endpoints.
The same statement holds for a time-integrated squared residual after
replacing $M_n(t)$ by their positive integrals.
\end{proposition}

The proof is elementary and useful operationally. A zero in mode one
forces $\omega=\pm a$ and $c=2a\omega$. Mode two excludes $\omega=-a$.
The leading quartic coefficient is $\sum n^4M_n>0$, while $c_*$ is
affine in $\omega$. A local optimizer can converge to a nonzero residual
minimum; the finite global search resolves this failure within this
one-parameter correction family. Floating-point root enumeration is
validated against the exact sum-of-squares identity, not presented as
an interval-arithmetic proof.

\subsection{An explicit distribution coupling and a particle limit}

\begin{theorem}[Phase coupling for the constructed field]
\label{thm:harmonic-coupling}
For $u_a$, a uniform grid $h=T/N$, and Euler initialized at $F_0(Z)$,
there is a constant $C$ independent of $N$ such that
\begin{equation}
 W_2(\widehat p_N,p_T)
 \le \|\widehat X_N-F_T(e^{-ia^2Th}Z)\|_{L^2(\rho)}
 \le Ch^2.
 \label{eq:harmonic-coupling}
\end{equation}
Meanwhile the common-start particle error has the strictly positive limit
\begin{equation}
 \lim_{h\downarrow0}\frac{\|\widehat X_N-F_T(Z)\|_{L^2}}h
 =a^2T\|DF_T(Z)JZ\|_{L^2(\rho)}>0.
 \label{eq:harmonic-particle-limit}
\end{equation}
\end{theorem}

The comparison curve $F_t(e^{-ia^2ht}Z)$ solves the explicit modified
field $u_a-(h/2)\acc_{u_a}$, which generates the same marginals.
A second-order Taylor expansion cancels Euler's quadratic local defect,
giving the strong $O(h^2)$ comparison in
Eq.~\eqref{eq:harmonic-coupling}; no infinite modified-equation series is
used. This gives a computable coupling bound for a genuinely nonradial
law. It is an upper bound on full $W_2$, not its exact value. The theorem
does not claim a positive $h^2$ limit for that optimal distance.
